\section{Introduction}
\label{sec:introduction}

Let $(X,\mathcal F,\mu)$ be a probability space and let $T:X\to X$
be invertible, bimeasurable, and measure preserving. We say that $T$
is \emph{mixing} if
\begin{equation}\label{eq:mixing}
 \mu(A\cap T^{-r}B)\longrightarrow\mu(A)\mu(B)
 \qquad (|r|\to\infty)
\end{equation}
for every $A,B\in\mathcal F$. For a fixed tuple of bounded measurable
functions, the consecutive multiple ergodic averages are
\[
 A_N(f_1,\ldots,f_n)(x)
 =\frac1N\sum_{k=1}^N\prod_{j=1}^n f_j(T^{jk}x).
\]
The problem is to determine their almost-everywhere behavior, not merely
their limit in norm. We prove the following result.

\begin{theorem}\label{thm:main}
Suppose that $T$ satisfies \eqref{eq:mixing}. For every integer
$n\ge2$ and every fixed $f_1,\ldots,f_n\in L^\infty(\mu)$,
\[
 A_N(f_1,\ldots,f_n)(x)
 \longrightarrow\prod_{j=1}^n\int_X f_j\dd\mu
 \qquad\text{for $\mu$-almost every }x,
\]
as $N$ tends to infinity through all positive integers. The probability
space need not be standard, and no rate of mixing is required.
\end{theorem}

The exceptional null set may depend on the system and the fixed tuple.
The theorem concerns the times $k,2k,\ldots,nk$ and does not require
a common exceptional set for all bounded functions.

\subsection{Context and prior work}

Arithmetic-progression averages became central to ergodic Ramsey theory
through Furstenberg's proof of Szemer\'edi's theorem
\cite{Furstenberg1977}. The distinction between norm and pointwise
convergence is substantial. Host and Kra proved $L^2$ convergence of
these averages for every finite length on arbitrary
probability-preserving systems \cite[Theorem~1.1]{HostKra2005}.
Ziegler obtained another proof through universal characteristic
factors \cite[Corollary~1.8]{Ziegler2007}.
For pointwise convergence, Bourgain's double recurrence theorem treats
two factors \cite{Bourgain1990}. Higher-length results have been
obtained under additional structural assumptions: K-systems
\cite[Theorem~4.2]{DerrienLesigne1996}, weakly mixing systems with
singular spectrum on their Pinsker factor \cite{Assani1998}, and
weakly mixing systems whose pairwise independent self-joinings are
independent \cite[Theorem~3.4]{GutmanHuangShaoYe2018}. Huang, Shao,
and Ye proved pointwise convergence for every finite length on
ergodic measure-distal systems \cite[Theorem~C]{HuangShaoYe2019}.

The present proof extends the fourfold mixing theorem of
\cite[Theorem~1.1]{FourfoldCompanion}. It uses two substantial
companion results. The first proves that ordinary mixing implies
mixing of every finite order, the assertion of Rokhlin's
multiple-mixing problem \cite[Theorem~1.1]{MixingCompanion}.
The second gives an oscillation estimate for triple averages at
distinct integer slopes on arbitrary invertible probability-preserving
systems \cite[Proposition~3.5]{SlopesCompanion}; together with its
profile approximation, this yields pointwise triple convergence
at distinct positive integer slopes without a mixing assumption.
Section~\ref{sec:inputs} states their exact forms and derives the
required triple pointwise theorem along all lengths. These companion
proofs are not reproduced here. The selected-law construction,
arbitrary-length extension argument, and tensor estimate used below
are proved in full.

\subsection{The arbitrary-length step}

We briefly describe what must change when the number of factors
increases. A failure of pointwise convergence allows an averaging
length to be chosen measurably at each starting point so that a
centered average stays positive. After passing to a suspension flow,
these choices define a probability law on $n$ output coordinates.
The law has the correct individual marginals and is invariant under
different speeds in the different coordinates. The same construction
works on extensions of the flow.

The method of \cite[Sections~3--6]{FourfoldCompanion} organizes
additional outputs into three groups. Their spatial coordinates are
independent, and the original spatial coordinate is independent of
any two groups; independence from all three groups is not imposed.
We retain this three-group structure even when there are $n$ slots.
Higher-order mixing and a tensor-packing estimate show that these
independence properties survive the selected joining construction.

Two changes make this approach work for arbitrary $n$. First, we
minimize the number of slots supporting a nonzero centered correlation.
This gives independence for every smaller collection of the specific
spatial lists being tested. A slot may read both the original spatial
coordinate and the outputs in any proper subset of the three groups.
That flexibility permits the outputs to be regrouped without imposing
independence on arbitrary collections of whole states.

Second, an extension is chosen so that conditioning on all distal
parameters together with a single state cannot improve under any
further extension. This is a conditional-expectation energy
construction, related to the sated-extension method
\cite{Austin2015} and adapted from
\cite[Section~6]{FourfoldCompanion}. Applying the joining construction
again gives a square array. The single-state conditioning identity
controls its distinguished row and column. Their parameter laws and
a Hilbert-space argument then produce a new nonzero correlation with
the same number of slots but one more output-only slot. Maximizing
that number gives the contradiction. Only the pointwise theorem for
at most three slots on arbitrary systems is used.

Section~\ref{sec:inputs} records the dynamical inputs.
Section~\ref{sec:selected} constructs the selected laws and their
conditional marginals. Section~\ref{sec:channels} proves preservation
of the three-group structure. Sections~\ref{sec:correlations}
and~\ref{sec:square} carry out the minimal-correlation and square
arguments. Appendix~\ref{sec:packing} proves the tensor estimate.
